\usetikzlibrary{arrows.meta, shapes.geometric, positioning}
\tikzset{
    umlactor/.style={
        rectangle, minimum width=1cm, minimum height=1.5cm,
        path picture={
            \draw[thick] (path picture bounding box.north) ++(0,-0.2) circle(0.2);
            \draw[thick] (path picture bounding box.north) ++(0,-0.4) -- ++(0,-0.5) coordinate(b);
            \draw[thick] (b) -- ++(-0.3,-0.5);
            \draw[thick] (b) -- ++(0.3,-0.5);
            \draw[thick] (path picture bounding box.north) ++(0,-0.6) -- ++(-0.4,0);
            \draw[thick] (path picture bounding box.north) ++(0,-0.6) -- ++(0.4,0);
        }
    },
    umlusecase/.style={
        ellipse, draw, thick, minimum width=3.5cm, minimum height=1.2cm,
        align=center, text width=3cm, font=\footnotesize
    },
    umlassoc/.style={thick, -Stealth},
    umlgen/.style={thick, -{Triangle[open, scale=1.5]}},
    umldep/.style={thick, dashed, -Stealth}
}

\section*{Цель}
Получение базовых навыков описания вариантов использования (Use
Cases, прецедентов) для подготовки тестирования.

\section*{Задачи}
\begin{enumerate}
    \item Изучить особенности построения UML-диаграммы автомата.
    \item Описать варианты использования в соответствии шаблоном, разработанным сообществом бизнес-аналитиков в России на основе BABOK Guide.
    \item Зафиксировать результаты в отчете.
\end{enumerate}

\section{Описание назначения техники тестирования на основе вариантов использования}
Тестирование на основе вариантов использования представляет собой метод проектирования тестов методом черного ящика. Основное назначение данной техники заключается в проверке поведения системы с точки зрения конечного пользователя и валидации бизнес-требований. Варианты использования позволяют выявить сценарии взаимодействия акторов с системой для достижения конкретных измеримых целей. На основе нормальных потоков, альтернативных путей и исключений, задокументированных в прецедентах, разрабатываются сценарии позитивного и негативного тестирования, обеспечивающие функциональное покрытие внешних требований к системе.

\section{Полная диаграмма вариантов использования}
На \cref{fig:main_diagram} представлена полная UML-диаграмма прецедентов для предметной области офиса продаж.

\begin{figure}[H]
    \centering
    \begin{tikzpicture}
        % Акторы
        \node[umlactor, label=below:{\footnotesize Пользователь}] (User) at (0, 0) {};
        \node[umlactor, label=below:{\footnotesize Кассир}] (Cashier) at (-2, 4) {};
        \node[umlactor, label=below:{\footnotesize Менеджер по продажам}, align=center] (Manager) at (2, 4) {};

        % Наследование (Generalization)
        \draw[umlgen] (Cashier) -- (User);
        \draw[umlgen] (Manager) -- (User);

        % Прецеденты
        \node[umlusecase] (Auth) at (3, -1) {Аутентификация};
        \node[umlusecase] (Pay) at (-4, 7) {Принять оплату};
        \node[umlusecase] (Receipt) at (0, 7) {Распечатать чек};
        \node[umlusecase] (Request) at (3.5, 7) {Оформить заявку};
        \node[umlusecase] (Increase) at (7.5, 7) {Увеличить количество\\товара на складе};
        \node[umlusecase] (Order) at (5.5, 4.5) {Заказать товар};
        \node[umlusecase] (Report) at (5.5, 1.5) {Построить отчет};
        \node[umlusecase] (Sales) at (9.5, 3) {Посмотреть отчет о\\продажах};
        \node[umlusecase] (Profit) at (9.5, 0) {Посмотреть отчет о\\прибыли};

        % Ассоциации акторов и прецедентов
        \draw[umlassoc] (User) -- (Auth);
        \draw[umlassoc] (Cashier) -- (Pay);
        \draw[umlassoc] (Cashier) -- (Receipt);
        \draw[umlassoc] (Manager) -- (Request);
        \draw[umlassoc] (Manager) -- (Order);
        \draw[umlassoc] (Manager) -- (Report);

        % Зависимости (Include / Extend)
        \draw[umldep] (Order) -- node[fill=white, inner sep=1pt, font=\scriptsize] {<<include>>} (Increase);
        \draw[umldep] (Sales) -- node[fill=white, inner sep=1pt, font=\scriptsize] {<<extend>>} (Report);
        \draw[umldep] (Profit) -- node[fill=white, inner sep=1pt, font=\scriptsize] {<<extend>>} (Report);
    \end{tikzpicture}
    \caption{Диаграмма прецедентов для офиса продаж}
    \label{fig:main_diagram}
\end{figure}

\section{Описание вариантов использования}
В соответствии с заданием выбраны два варианта использования для разных акторов: \textit{Принять оплату} для актора \textit{Кассир} и \textit{Заказать товар} для актора \textit{Менеджер по продажам}. Подробное описание представлено в \cref{tab:uc1} и \cref{tab:uc2}.

\begin{longtblr}[
    caption = {Описание варианта использования \textit{Принять оплату}},
    label = {tab:uc1},
]{
    colspec = { X[1,l,m] X[2,l,m] },
    hlines, vlines,
}
    ID Варианта использования: & UC1 \\
    Наименование варианта использования: & Принять оплату \\
    Кем создан: Студент & Кем в последний раз изменен: Студент \\
    Дата создания: 01.11.2024 & Дата последнего изменения: 01.11.2024 \\
    Акторы: & Кассир \\
    Описание: & Основной поток событий: Кассир принимает оплату от клиента за выбранные товары для успешного завершения процесса покупки. \\
    Предварительные условия: & Кассир авторизован в системе, товары добавлены в чек, сформирована итоговая сумма к оплате. \\
    Постусловие: & 1. Система регистрирует оплату в базе данных. \par 2. Статус заказа изменяется на \textit{Оплачен}. \\
    Нормальный ход событий: & 
    1. Прецедент начинается, когда кассир инициирует процесс оплаты. \par
    2. Система запрашивает выбор способа оплаты. \par
    3. Кассир выбирает безналичный расчет. \par
    4. Система ожидает подтверждения от банковского терминала. \par
    5. Терминал подтверждает успешную транзакцию. \par
    6. Система выводит сообщение об успехе и переходит к формированию чека. \\
    Альтернативный ход событий: & 
    3a. Кассир выбирает оплату наличными. \par
    4a. Кассир вводит полученную от клиента сумму. \par
    5a. Система рассчитывает и отображает сдачу. \par
    6a. Кассир подтверждает завершение расчета. \\
    Исключения: & 
    5b. Банковский терминал отклоняет карту (недостаточно средств). Система выводит сообщение об ошибке и просит выбрать другой способ оплаты. \par
    5c. Отсутствует связь с терминалом. Система прерывает операцию и просит повторить попытку. \\
    Содержит: & \\
    Приоритет: & Высший \\
    Частота использования: & Многократно в течение рабочей смены \\
    Бизнес-правила: & Оплата должна строго соответствовать фискальным требованиям. \\
    Специальные требования: & Интеграция с эквайринговым терминалом. \\
    Предпосылки (предположения): & Клиент имеет при себе средства для оплаты заказа. \\
    Примечания и вопросы: & \\
\end{longtblr}

Графическое представление прецедента \textit{Принять оплату} показано на \cref{fig:uc1_graph}.

\begin{figure}[H]
    \centering
    \begin{tikzpicture}
        \node[umlactor, label=below:{\footnotesize Кассир}] (Cashier) at (0, 0) {};
        \node[umlusecase] (Pay) at (4, 0) {Принять оплату};
        
        \draw[umlassoc] (Cashier) -- (Pay);
    \end{tikzpicture}
    \caption{Графическое представление варианта использования \textit{Принять оплату}}
    \label{fig:uc1_graph}
\end{figure}

\begin{longtblr}[
    caption = {Описание варианта использования \textit{Заказать товар}},
    label = {tab:uc2},
]{
    colspec = { X[1,l,m] X[2,l,m] },
    hlines, vlines,
}
    ID Варианта использования: & UC2 \\
    Наименование варианта использования: & Заказать товар \\
    Кем создан: Студент & Кем в последний раз изменен: Студент \\
    Дата создания: 01.11.2024 & Дата последнего изменения: 01.11.2024 \\
    Акторы: & Менеджер по продажам \\
    Описание: & Основной поток событий: Менеджер оформляет заказ на поставку отсутствующего товара для пополнения запасов на складе. \\
    Предварительные условия: & Менеджер авторизован в системе, выявлен дефицит определенной позиции товара. \\
    Постусловие: & 1. Заявка на товар отправлена поставщику. \par 2. Система обновляет статус ожидаемого товара. \\
    Нормальный ход событий: & 
    1. Прецедент начинается, когда менеджер выбирает опцию \textit{Заказать товар}. \par
    2. Система отображает форму заказа. \par
    3. Менеджер вводит артикул товара и необходимое количество. \par
    4. Система проверяет корректность данных. \par
    5. Система вызывает прецедент \textit{Увеличить количество товара на складе}. \par
    6. Система сохраняет заказ и выводит уведомление об успешном оформлении. \\
    Альтернативный ход событий: & 
    3a. Менеджер выбирает товар из автоматически сформированного списка дефицитных позиций вместо ручного ввода артикула. \par
    4a. Менеджер отменяет создание заказа, нажимая кнопку \textit{Отмена}. \\
    Исключения: & 
    4b. Введен некорректный артикул. Система подсвечивает поле и просит исправить ввод. \par
    4c. Запрошенное количество превышает допустимый лимит. Система выдает предупреждение и блокирует отправку заявки до подтверждения руководством. \\
    Содержит: & <<include>> Увеличить количество товара на складе \\
    Приоритет: & Средний \\
    Частота использования: & Несколько раз в неделю \\
    Бизнес-правила: & Заказ формируется только у утвержденных поставщиков. \\
    Специальные требования: & Доступ к актуальному каталогу товаров. \\
    Предпосылки (предположения): & Поставщик готов принять заявку. \\
    Примечания и вопросы: & \\
\end{longtblr}

Графическое представление прецедента \textit{Заказать товар} показано на \cref{fig:uc2_graph}.

\begin{figure}[H]
    \centering
    \begin{tikzpicture}
        \node[umlactor, label=below:{\footnotesize Менеджер по продажам}, align=center] (Manager) at (0, 0) {};
        \node[umlusecase] (Order) at (4, 0) {Заказать товар};
        \node[umlusecase] (Increase) at (9.5, 0) {Увеличить количество\\товара на складе};
        
        \draw[umlassoc] (Manager) -- (Order);
        \draw[umldep] (Order) -- node[above, font=\scriptsize] {<<include>>} (Increase);
    \end{tikzpicture}
    \caption{Графическое представление варианта использования \textit{Заказать товар}}
    \label{fig:uc2_graph}
\end{figure}

\section*{Выводы}
В ходе выполнения практической работы была достигнута поставленная цель — получены базовые навыки описания вариантов использования для подготовки тестирования. Были изучены особенности построения UML-диаграммы прецедентов и спецификации требований. На основе предложенной диаграммы были выделены и подробно описаны два варианта использования для разных акторов (Кассир и Менеджер по продажам) с применением шаблона, основанного на BABOK Guide. В процессе работы были проанализированы нормальные и альтернативные ходы событий, а также исключения, что является важным этапом для последующей разработки позитивных и негативных тест-кейсов. Возникшая сложность при определении альтернативных потоков была решена путем детальной декомпозиции возможных действий пользователя и реакций системы. Полученные навыки позволят в дальнейшем эффективно проектировать тестовые сценарии на основе задокументированных требований.